\part{Mathematical Foundations and Economic Results}
\label{part:IX}

Part IX consolidates the formal machinery and states the principal results with hypotheses, organized by theorem family. It is the formal centre of the book.

\begin{housing}The housing model is stated formally, its viability kernel is computed, and the finite-state companion reproduces the results and their counterexamples.\end{housing}

\chapter{Transition Systems and History Compression}\label{ch:41}
Consolidates the formal machinery used throughout the book. Establishes sufficient-state conditions, observational equivalence, and history-dependent constraints, with careful treatment of what a chosen state representation forgets. The representation family is stated here as a single factorization lemma with a dynamic congruence version.
\chspec{\Pref{P01}, \Pref{P02}, \Pref{P11}, \Pref{P19}.}
\scaffold

\chapter{Orders, Dominance, and Incomparability}\label{ch:42}
Equips continuation profiles with partial orders reflecting specified dimensions of improvement. Establishes conditions for dominance and explains why inclusion, cardinality, and scalar utility are different comparisons. Any impossibility result concerning scalarization states the preservation requirements that generate it, and the multi-utility companion is given.
\chspec{\Pref{P13}.}
\scaffold

\chapter{Reachability, Robustness, and Recovery}\label{ch:43}
Develops finite-state algorithms and selected continuous models for viability and corrective reachability, with the fixed-point theorem and comparative statics. Compares worst-case guarantees with probabilistic guarantees and analyzes the effects of model uncertainty, including almost-sure and risk-budget variants. Includes the waiting theorem.
\chspec{\Pref{P03}, \Pref{P04}, \Pref{P05}, \Pref{P06}, \Pref{P07}, \Pref{P08}, \Pref{P14}, \Pref{P23}, \Pref{P25}, \Pref{P26}.}
\scaffold

\chapter{Composition and Systemic Failure}\label{ch:44}
Assembles local contracts into networks with shared resources and timing constraints. Counterexamples show how individually feasible commitments become jointly infeasible. Positive results identify sufficient compatibility and resource conditions, including independent composition, scarcity prices, and circular assume-guarantee reasoning.
\chspec{\Pref{P09}, \Pref{P10}, \Pref{P18}, \Pref{P24}, \Pref{P27}.}
\draftnote{Candidate case: a transactional outbox records the obligation to send an event in the same transaction as the order, which prevents one class of omission. Without further assumptions it does not guarantee eventual delivery or prevent duplicate effects at the recipient, which separates accounting consistency from operational progress (P20, P27).}
\scaffold

\chapter{Measurement and Empirical Identification}\label{ch:45}
Translates theoretical quantities into observable indicators and study designs. Maintenance records, switching costs, recovery times, service interruptions, and obligation fulfillment provide evidence. Explains which causal claims require experiments, natural experiments, longitudinal data, or structural assumptions.
\chspec{\Pref{P12}, \Pref{P15}, \Pref{P16}, \Pref{P17}, \Pref{P28}, \Pref{P29}.}
\scaffold
